% Documentation of the kpi-thesis bundle, typeset with the class itself (lualatex).
\documentclass{kpi-thesis}
\title{kpi-thesis}
% A short manual, not a thesis: sections run on, tables are numbered 1, 2, ..., code is larger.
\RemoveFromHook{cmd/section/before}[kpi-thesis]
\renewcommand\thetable{\arabic{table}}
\fvset{fontsize=\setstretch{1}\fontsize{11}{13}\selectfont}
\begin{document}

\begin{center}
  \textbf{\large kpi-thesis}\par
  Клас LaTeX для дипломного проєкту бакалавра й курсових робіт\par
  КПІ ім.~Ігоря Сікорського, ФІОТ\par
  \medskip
  Версія 2.0.0, 17 листопада 2026 р.\par
  Олександр Соколов, \url{https://github.com/sokolovgit/kpi-thesis}
\end{center}

\bigskip
\noindent\textit{The kpi-thesis class typesets a bachelor diploma project and courseworks of
Igor Sikorsky Kyiv Polytechnic Institute (Faculty of Informatics and Computer Engineering)
following the Ukrainian standard DSTU~3008:2015 and the department guide: title page,
assignment, abstracts, project sheet with ESKD title blocks, A3 drawing sheets, ДСТУ tables,
formula explanations and lists. The documentation is in Ukrainian, the language of the
documents it produces.}

\structural{ЗМІСТ ПАКЕТА}

Пакет складається з класу і пакетів, які клас підключає сам або які додає документ:
\begin{itemize}
  \item \texttt{kpi-thesis.cls} --- поля, шрифти, заголовки, підписи рисунків і таблиць,
    додатки, список джерел, нумерація сторінок;
  \item \texttt{kpi-thesis-titles.sty} --- дані роботи (\verb|\kpiset|), титульний лист,
    завдання, анотації, відомість проєкту, титули документів;
  \item \texttt{kpi-thesis-graphics.sty} --- аркуші графічного матеріалу А3 зі штампом ЄСКД;
  \item \texttt{kpi-thesis-tables.sty}, \texttt{kpi-thesis-where.sty},
    \texttt{kpi-thesis-lists.sty}, \texttt{kpi-thesis-stats.sty} --- таблиці з
    «Продовження таблиці», пояснення до формул, переліки «а)», підрахунок сторінок, рисунків,
    таблиць, джерел і додатків для анотації. Клас підключає їх сам; кожен можна взяти й
    окремо.
\end{itemize}

Шаблони диплома й курсової, лінтер нормоконтролю \texttt{just check} і плагін для Claude Code
живуть у репозиторії проєкту; там само --- покрокове встановлення для macOS, Linux і Windows.

\structural{ЗБИРАННЯ}

Рушій --- LuaLaTeX (\texttt{fontspec}). XeLaTeX ще працює, але клас попереджає: команда LaTeX
більше не розвиває XeTeX. pdfLaTeX не підтримується.

Методичка вимагає шрифти Times New Roman і Courier New. Якщо їх немає в системі (Overleaf,
Linux без \texttt{ttf-mscorefonts-installer}), клас бере метрично сумісні TeX Gyre Termes і
TeX Gyre Cursor і пише попередження. Для здачі потрібен саме Times New Roman.

\structural{ОПЦІЇ КЛАСУ}

\begin{Verbatim}
\documentclass[coursework, main-font={Times New Roman}]{kpi-thesis}
\end{Verbatim}

\begin{kpitable}[0]{|L{4cm}|L{11.3cm}|}{Опції класу}{\textbf{Опція} & \textbf{Що робить}}
\texttt{diploma} & документ диплома (типово) \\ \hline
\texttt{coursework} & курсова робота \\ \hline
\texttt{main-font=} & основний шрифт замість Times New Roman \\ \hline
\texttt{mono-font=} & моноширинний шрифт замість Courier New \\ \hline
інші & передаються класу \texttt{extarticle} \\ \hline
\end{kpitable}

Там, де методичка курсової розходиться з дипломною, у преамбулі ставиться
\verb|\kpisetup{...}|:

\begin{kpitable}[0]{|L{6.2cm}|L{9.1cm}|}{Ключі \texttt{\textbackslash kpisetup}}{\textbf{Ключ} & \textbf{Що робить}}
\texttt{left-margin=20mm} & ліве поле (типово 30 мм) \\ \hline
\texttt{hyphenation=false} & без переносів \\ \hline
\texttt{enumerate=arabic} & перелік «1)», вкладений «а)» (типово навпаки) \\ \hline
\texttt{table-caption=left} & підпис таблиці від лівого краю, а не з абзацу \\ \hline
\texttt{appendix-title=inline} & «Додаток А. Назва» одним рядком \\ \hline
\texttt{url-breaks=anywhere} & URL рветься будь-де (\texttt{xurl}) \\ \hline
\texttt{algorithms} & \texttt{algorithm} і \texttt{algpseudocode}, «Алгоритм 1.1 -- Назва» \\ \hline
\end{kpitable}

\structural{КОМАНДИ}

\begin{kpitable}[0]{|L{6.2cm}|L{9.1cm}|}{Команди класу}{\textbf{Команда} & \textbf{Призначення}}
\verb|\structural{ВСТУП}| & структурний елемент без номера, у змісті \\ \hline
\verb|\mainmatter| & нумерація сторінок з наступної після титулу \\ \hline
\texttt{kpitable} & таблиця з «Продовження таблиці N» і повтором шапки \\ \hline
\texttt{where} & пояснення до формули: «де» і символи з тире \\ \hline
\verb|\appendices|, \verb|\appendix{Назва}| & додатки «Додаток А» \\ \hline
\verb|\appendixpdf| & готовий PDF (звіт подібності тощо) як додаток \\ \hline
\verb|\codefile|, \verb|\codeinput| & текст програми (\texttt{fvextra}, 8~пт) \\ \hline
\verb|\diagramPage|, \verb|\graphicsSheet| & аркуш А3 зі штампом \\ \hline
\verb|\todo{...}| & червона заглушка, яку рахує \texttt{just check} \\ \hline
\end{kpitable}

Таблиця з продовженням на наступних сторінках:
\begin{Verbatim}
\begin{kpitable}{|L{4cm}|L{10cm}|}{Назва\label{tab:key}}
                {Поле & Опис}
  Рядок & Значення \\ \hline
\end{kpitable}
\end{Verbatim}

Пояснення символів після формули --- без двокрапки після «де», кожен символ з нового рядка:
\begin{Verbatim}
\begin{where} M & маса, кг; \\ F & сила, Н. \end{where}
\end{Verbatim}

\structural{ДАНІ РОБОТИ}

Тема, ПІБ, група, шифри й дати задаються один раз через \verb|\kpiset{...}| з пакета
\texttt{kpi-thesis-titles}; титул, завдання, штампи й колонтитули беруть їх звідти.
Курсова:
\begin{Verbatim}
\documentclass[coursework]{kpi-thesis}
\usepackage{kpi-thesis-titles}
\kpiset{title = {Тема}, discipline = {Дисципліна}, group = {ІП-XX},
        student = {Прізвище Ім'я По батькові},
        advisor = {...}, year = {2027}}
\begin{document}
\kpiCourseworkTitle
\mainmatter
\tableofcontents
\end{document}
\end{Verbatim}

Диплом складається з кількох документів (титул, ТЗ, пояснювальна записка, текст програми,
ПМТ, керівництво користувача, графічний матеріал) зі спільним \texttt{meta.tex}; повний шаблон
і рецепти збирання --- у репозиторії.

\structural{ЛІЦЕНЗІЯ}

MIT. Вимоги до оформлення: ДСТУ~3008:2015, ДСТУ~8302:2015 і методичка кафедри інформатики та
програмної інженерії; методичку й вимоги час від часу оновлюють, тож звіряйтеся з актуальною
редакцією і з керівником.

\end{document}
